% LLM Lifecycle Lab 教程全景图：数据下载 -> 治理 -> 训练就绪数据 -> 下游阶段
\documentclass[border=24pt]{standalone}
\usepackage[fontset=none]{ctex}
\setCJKmainfont{Noto Serif CJK SC}
\setCJKsansfont{Noto Sans CJK SC}
\setmainfont{Noto Sans CJK SC}
\usepackage{tikz}
\usetikzlibrary{arrows.meta, positioning, fit, calc}

\definecolor{cblue}{HTML}{2563EB}
\definecolor{cgreen}{HTML}{059669}
\definecolor{cpurple}{HTML}{7C3AED}
\definecolor{corange}{HTML}{D97706}
\definecolor{cred}{HTML}{DC2626}
\definecolor{cgray}{HTML}{6B7280}
\definecolor{bgblue}{HTML}{EFF6FF}
\definecolor{bggreen}{HTML}{ECFDF5}
\definecolor{bgpurple}{HTML}{F5F3FF}
\definecolor{bgorange}{HTML}{FFF7ED}
\definecolor{bgred}{HTML}{FEF2F2}
\definecolor{bggray}{HTML}{F3F4F6}

\tikzset{
  block/.style 2 args={draw=#1, fill=#2, rounded corners=3pt, inner sep=5pt, font=\small},
  step/.style={draw=#1, fill=white, rounded corners=2pt, inner sep=3pt, font=\footnotesize, text=#1, align=center},
  box/.style 2 args={draw=#1!50, fill=#2, rounded corners=5pt, inner sep=6pt},
}

\begin{document}
\begin{tikzpicture}[node distance=8pt, >=Stealth, thick]

% ===== 左：原始数据 =====
\node[box={cred}{bgred}, align=center] (raw) at (0,0) {\textbf{原始数据} \texttt{data/raw}\\[2pt]
{\footnotesize 只读，不修改}};

\node[block={cred}{bggray}, align=left, below=10pt of raw.south] (raw-train) {\footnotesize\textbf{训练/微调源}\\
TinyStories、Wikitext-103\\
tigerbot-law、alpaca-gpt4-zh\\
ultrafeedback、gsm8k、ChartQA};
\node[block={cred}{bggray}, align=left, below=6pt of raw-train.south] (raw-eval) {\footnotesize\textbf{仅评测源}\\
HellaSwag、C-Eval};

% ===== 中：治理管线 =====
\node[box={cblue}{bgblue}, align=center] (gov) at (7.0,2.0) {\textbf{数据治理管线}\\[2pt]
{\footnotesize 统一流程，逐数据集执行}};

% 5x2 步骤网格：居中于 gov 下方，两行独立 y 坐标，列距 1.9cm，行距 0.8cm
\foreach \s/\n [count=\i] in {
  盘点/1, 许可证/2, schema 检查/3, 有界抽样/4, 清洗/5}{
  \pgfmathsetmacro{\colx}{(\i-2.5)*1.9}
  \node[step=cblue, minimum width=1.6cm, anchor=center] at ([xshift=\colx cm, yshift=-1.35cm]gov.center) {\n.\ \s};
}
\foreach \s/\n [count=\i] in {
  去重/6, 分组切分/7, Parquet/8, token 统计/9, manifest/10}{
  \pgfmathsetmacro{\colx}{(\i-2.5)*1.9}
  \node[step=cblue, minimum width=1.6cm, anchor=center] at ([xshift=\colx cm, yshift=-2.15cm]gov.center) {\n.\ \s};
}

% ===== 右：产物 =====
\node[box={cgreen}{bggreen}, align=center] (out) at (16.2,1.4) {\textbf{训练就绪数据}\\[2pt]
{\footnotesize \texttt{data/processed} + \texttt{data/manifests}}};

\node[block={cgreen}{bggray}, align=left, below=10pt of out.south] (out-p) {\footnotesize\textbf{Parquet 分片}\\
train / validation / test\\
(或 rm\_train / dpo\_train ...)};
\node[block={cgreen}{bggray}, align=left, below=6pt of out-p.south] (out-m) {\footnotesize\textbf{Manifest 清单}\\
许可证、revision、seed、\\
行数、token 数、切分策略};

% ===== 下：下游阶段 =====
\node[box={cpurple}{bgpurple}, align=center] (down) at (7.2,-2.9) {\textbf{下游阶段（后续教程）}\\[2pt]
{\footnotesize 预训练 $\rightarrow$ CPT $\rightarrow$ SFT $\rightarrow$ RM/DPO $\rightarrow$ GRPO $\rightarrow$ VLM $\rightarrow$ 评测 $\rightarrow$ 量化 $\rightarrow$ 部署}};

% ===== 箭头（折线：先垂直后水平，避免穿过下方步骤网格） =====
\draw[cred, ->, thick] (raw-train.east) |- node[above, font=\footnotesize]{治理} (gov.west);
\draw[cred, ->, thick, dashed] (raw-eval.east) |- node[above, font=\footnotesize]{仅记录} ([yshift=-8pt]gov.west);
\draw[cblue, ->, thick] (gov.east) -- node[above, font=\footnotesize]{产物} (out.west);
\draw[cgreen, ->, thick] (out.south) -- node[right, font=\footnotesize]{消费} (down.east);
\draw[cblue, ->, thick, dashed] ([yshift=-30pt]gov.south) -- node[right, font=\footnotesize]{流水线可复用} (down.north);

\end{tikzpicture}
\end{document}
